\documentclass[11pt,a4paper,final,addpoints]{exam}
\usepackage[utf8]{inputenc}
\usepackage[portuguese]{babel}
\usepackage[T1]{fontenc}
\usepackage{amsmath}
\usepackage{amsfonts}
\usepackage{amssymb}
\usepackage{graphicx}
\usepackage{pstricks,pst-circ,pst-osci,pst-magneticfield}
\usepackage[left=1cm,right=1cm,top=2cm,bottom=2cm,twoside]{geometry}
\author{Tiago Á. Oliveira}
\title{Relatorio Individual}


\pagestyle{headandfoot}

\extraheadheight[1.5in]{-.25in}

\firstpageheader{\textbf{Prova Extraordin\'{a}ria}}
{\includegraphics[scale=1]{figs/logo_UA}
\linebreak\linebreak
\LARGE\textbf{Eletricidade \& Magnetismo}
\normalsize\linebreak
(2014-2015)
\linebreak\linebreak
\textbf{Relat\'{o}rio Individual de Laborat\'{o}rio}
\linebreak
17 de dezembro de 2014
\linebreak}
{\textbf{Vers\~{a}o \'{U}nica}}

\runningheader{\oddeven{}{Eletricidade \& Magnetismo (2014-15)}}{}{\oddeven{Relat\'{o}rio Individual de Laborat\'{o}rio}{}}
\headrule

\footrule
\footer{\oddeven{}{P\'{a}gina \thepage}}{\iflastpage{Fim!}}{\oddeven{P\'{a}gina \thepage}{}}

% RESPOSTAS 
\printanswers
\renewcommand{\solutiontitle}{\noindent\textbf{Resposta:}}


\begin{document}
\begin{center}
\makebox[\textwidth]{\textbf{Nome:}~\underline{\hspace{11cm}}~\textbf{N.\textsuperscript{o} Mecanogr\'{a}fico:}~\underline{\hspace{2.5cm}}}
\end{center}

\noindent\textbf{Recomenda\c{c}\~{o}es:}
\begin{itemize}
\item O relat\'{o}rio \'{e} realizado individualmente;

\item Leia atentamente todo o enunciado da prova e se tiver d\'{u}vidas coloque-as ao docente;

\item Tem a dura\c{c}\~{a}o exata de 2 horas, sem qualquer toler\^{a}ncia;

\item O enunciado do relat\'{o}rio tem apenas uma vers\~{a}o. 

\item \underline{Dever\~{a}o justificar todas as respostas};

\item N\~{a}o ser\'{a} permitida a consulta dos guias, fichas ou dos registos efetuados durante as aulas;

\item N\~{a}o haver\'{a} formul\'{a}rio para as leis fundamentais (Lei de Ohm, Kirschoff, Faraday, Lenz, etc.). Equa\c{c}\~{o}es de dedu\c{c}\~{a}o complexa/demorada, caso sejam necess\'{a}rias, estar\~{a}o presentes no enunciado;

\item Podem usar m\' {a}quina de calcular com capacidades gr\'{a}ficas;

\item O uso do computador/tablet est\'{a} restringido apenas para uso no tratamento de dados (Excel, MatLAB, etc.) n\~{a}o podendo ser usado para consulta de PDF's ou liga\c{c}\~{a}o \`{a} Internet;

\item O telem\'{o}vel dever\'{a} estar desligado durante a prova.
\end{itemize}

\noindent\textbf{Cota\c{c}\~{o}es:}
\begin{center}
\hqword{\textit{Quest\~{a}o}}
\hpword{\textit{Pontos}}
\hsword{\textit{Obtidos}}
\htword{\textbf{Total}}
\gradetable[h][questions]
\linebreak
\end{center}

\noindent\textbf{Constantes fundamentais:}
\begin{center}
\begin{tabular}{cc}
$\varepsilon_0 = 8.85\times10^{-12}~\text{C}^2/\text{Nm}^2$ & $\mu_0 = 4\pi\times10^{-7}~\text{N}/\text{A}^2$ \\ 
$e=1.60\times10^{-19}~\text{C}$ & $m_e=9.11\times10^{-31}~\text{kg}$ \\ 
\end{tabular} 
\end{center}


\boxedpoints
\pointsinmargin
\newpage
\begin{questions}
% Versao Unica
% Questao sobre o Eletrometro
% Pontuacao - 25 pontos
\question[25]
\textbf{Trabalho Pr\'{a}tico 1 : Eletroest\'{a}tica}

Explique o princ\'{i}pio de funcionamento do eletr\'{o}metro que usou na realiza\c{c}\~{a}o dos trabalhos pr\'{a}ticos 1 e 2 (Eletroest\'{a}tica e Condensador de Placas Paralelas), respondendo \`{a}s seguintes quest\~{o}es:
\begin{parts}
\part
Como \'{e} que este permite a medi\c{c}\~{a}o da carga?
\part
Porque \'{e} necessario manter a entrada do eletr\'{o}metro ligada \`{a} terra sempre que n\~{a}o se esteja a fazer medi\c{c}\~{o}es?
\end{parts}

\question[25]
\textbf{Trabalho Pr\'{a}tico 3 : Circuitos de Corrente Cont\'{i}nua}

Durante a realiza\c{c}\~{a}o do trabalho pr\'{a}tico 3 foram apresentadas duas configura\c{c}\~{o}es destintas de medi\c{c}\~{a}o de uma resit\^{e}ncia desconhecida $R_x$, a partir da medi\c{c}\~{a}o, atrav\'{e}s de um amper\'{i}metro, da corrente $I$ que a atravessa e da medi\c{c}\~{a}o da queda de tens\~{a}o nesta, atrav\'{e}s de um voltimetro, conforme ilustra a figura seguinte.

\begin{figure}[h]
\begin{center}
\begin{pspicture}[showgrid=false](8,6)
\pnodes(1,1){A}(1,3){B}(4,3){C}(7,3){D}(7,1){E}(1,5){F}(4,5){G}(7,5){H}
\vdc[labeloffset=-1.0](A)(E){\textbf{(a)}}
\resistor[dipolestyle=zigzag](C)(D){$R_x$}
\circledipole[labeloffset=0](B)(C){\Large\textbf{A}}
\circledipole[labeloffset=0](G)(H){\Large\textbf{V}}
\wire[arrows=-*](G)(C)
\wire(A)(B)
\wire(D)(E)
\wire(D)(H)
\end{pspicture}
\begin{pspicture}[showgrid=false](8,6)
\pnodes(1,1){A}(1,3){B}(4,3){C}(7,3){D}(7,1){E}(1,5){F}(4,5){G}(7,5){H}
\vdc[labeloffset=-1.0](A)(E){\textbf{(b)}}
\resistor[dipolestyle=zigzag](C)(D){$R_x$}
\circledipole[labeloffset=0](B)(C){\Large\textbf{A}}
\circledipole[labeloffset=0](F)(H){\Large\textbf{V}}
\wire(F)(B)
\wire(A)(B)
\wire(D)(E)
\wire(D)(H)
\end{pspicture}
\end{center}
\caption{\label{fig:curtalongaderiv}\textbf{(a)}~-~Curta deriva\c{c}\~{a}o;~\textbf{(b)}~-~Longa deriva\c{c}\~{a}o;.}
\end{figure}

Sabendo que $R_A$ e $R_V$ s\~{a}o as resist\^{e}ncias internas do amperimetro e volt\'{i}metro, respertivamente, se pretender-mos medir uma resist\^{e}ncia com $R_x \gg R_A$, qual das duas configura\c{c}\~{o}es anteriores deveremos escolher? \underline{Justifique convenientemente a sua resposta}.
\question
\textbf{Trabalho Pr\'{a}tico 6 : Indu\c{c}\~{a}o Eletromagn\'{e}tica}

Recorde o trabalho pr\'{a}tico sobre a Indu\c{c}\~{a}o eletromagn\'{e}tica.

\begin{parts}
\part[10]
Na primeira parte deste trabalho montou o circuito da Fig.~\ref{fig:bobinas}, com o objetivo de estudar a Lei de Faraday e de Lenz.

Descreva o que observou no movimento dos ponteiros do galvan\'{o}metro quando, ap\'{o}s ter mantido o circuito ligado, desligou a fonte de tens\~{a}o e, \`{a} luz das leis de Faraday e de Lenz, indique o sentido da corrente induzida nesse instante (em que desligou) no circuito da bobina 2.

\begin{figure}[h]
\centering
\begin{pspicture}[showgrid=false](9,5)
\pnodes(1,1){A}(1,4){B}(4,4){C}(4,1){D}(5,1){E}(5,4){F}(8,4){G}(8,1){H}
\vdc[labeloffset=1](A)(D){$15~\text{V}$}
\resistor[dipolestyle=zigzag,labelangle=:U](A)(B){$5~\Omega$}
\coil[dipolestyle=curved,
      labelangle=:U,
      intensitylabel=$I$,
      labeloffset=-.5](C)(D){$1200~\text{esp.}$}
\newSwitch[tensionlabel=Bobina 1,
           tensioncolor=white,
           ison=false,
           tensionlabeloffset=0.7](B)(C){}

\coil[dipolestyle=curved,
      labelangle=:U,
      labeloffset=-.5](E)(F){$3600~\text{esp.}$}
\wire(G)(H)      
\circledipole[tensionlabel=Bobina 2,
           tensioncolor=white,
           tensionlabeloffset=0.7,
           labeloffset=0.0](F)(G){\Large\textbf{G}}
\wire(E)(H)

\end{pspicture}
\caption{\label{fig:bobinas}}
\end{figure}

\part\label{part:m}
Considere um circuito no qual est\~{a}o ligados em s\'{e}rie um gerador de sinal, uma resist\^{e}ncia de $R\pm\Delta R=100\pm5~\Omega$ e um solenoide (daqui em diante referido como \emph{prim\'{a}rio}) de raio $r\pm\Delta r=1.38\pm0.01~\text{cm}$ e $N/l\pm\Delta N/l=3000\pm15$ espiras por metro. Uma bobina (daqui em diante referida como \emph{secund\'{a}rio}) de $N_b\pm\Delta N_b=1200\pm6$ espiras envolve o \emph{prim\'{a}rio}. Aplica-se um sinal triangular ao circuito \emph{prim\'{a}rio} tendo-se observado nos terminais da resist\^{e}ncia o sinal do Canal A da Fig.~\ref{fig:osci} (\emph{tracejado}). O sinal induzido no \emph{secund\'{a}rio} \'{e} o que se observa no Canal B da mesma figura (\emph{cheio}).

\begin{subparts}
\subpart[30]
Determine o coeficiente de indu\c{c}\~{a}o m\'{u}tua observado e compare-o com o esperado, sabendo que a for\c{c}a eletromotriz (f.e.m.) induzida no \emph{secund\'{a}rio} \'{e} dada pela Eq.~\ref{eq:fem}.

\begin{equation}
\label{eq:fem}
\varepsilon=\pi\mu_0\frac{N}{l}r^2N_b\frac{\mathrm d i}{\mathrm d t}
\end{equation}

\begin{figure}[h]
\begin{center}
\newpsstyle{Dash}{linestyle=dashed,
linecolor=black,linewidth=0.035,plotpoints=50}
\psscalebox{1}{
\Oscillo[Wave1=\TriangleA,
         Wave2=\RectangleB,
    amplitude1=4.00,
    amplitude2=-0.3,
       period1=1.2,
       period2=1.2,
        phase1=0.0,
        phase2=0.0,        
    sensivity1=1.0,
    sensivity2=0.2,
       timediv=0.15,
       offset2=+0.9,        
    plotstyle1=Dash,
      AllColor=false]}
\caption{\label{fig:osci}Visor do oscilosc\'{o}pio. O Canal A est\'{a} representado a tracejado enquanto o Canal B a cheio.}
\end{center}
\end{figure}

\subpart[15]
Justique convenientemente a forma do sinal induzido (Fig.~\ref{fig:osci}~-~\textit{cheio}).

\subpart[10]
Porque \'{e} que o sinal induzido (Fig.~\ref{fig:osci}~-~\textit{cheio}) apresenta flutua\c{c}\~{o}es nas extremidades?

\subpart[10]
Se o sinal aplicado (Canal A) fosse do tipo "dente-de-serra", o sinal induzido (Canal B) seria triangular?

\end{subparts}
\end{parts}
\newpage
\question
\textbf{Trabalho Pr\'{a}tico 5 : Bobinas de Helmholtz}

Recorde o trabalho pr\'{a}tico sobre o estudo do campo magn\'{e}tico produzido pelas Bobinas de Helmholtz.

\begin{parts}
\part[20]
No decorrer deste trabalho usou as bobinas de Helmholtz em diferentes configura\c{c}\~{o}es: com corrente a fluir apenas numa bobina; com corrente a fluir no mesmo sentido nas duas bobinas; e com corrente a fluir nas duas bobinas em sentidos opostos. Na Fig.~\ref{fig:helmholtzfield} encontra-se as linhas de campo produzidas por uma dessas configura\c{c}\~{o}es. Indique a qual configura\c{c}\~{a}o corresponde as linhas de campo ilustradas, justificando convenientemente a sua escolha.

\begin{figure}[ht]
\centering
\psset{unit=1,AntiHelmholtz,N=2,
R=2,pointsB=1000,pointsS=1000,PasB=0.001,PasS=0.003,nS=10,
nL=2,drawSelf=true,styleSpire=styleSpire,styleCourant=sensCourant}
\newpsstyle{grille}{subgriddiv=0,gridcolor=blue!50,griddots=10}
\newpsstyle{cadre}{linecolor=yellow!50}
\newpsstyle{styleSpire}{linecap=1,linecolor=red,linewidth=4\pslinewidth}
\newpsstyle{sensCourant}{linecolor=red,linewidth=1\pslinewidth}
\begin{pspicture*}[showgrid](-4,-4)(4,4)
\rput{90}{\psframe*[linecolor={[HTML]{996666}}](-7,6)(7,6)
\psmagneticfield[linecolor=black]}
\end{pspicture*}
\caption{\label{fig:helmholtzfield}Linhas de campo magn\'{e}tico produzido pelas Bobinas de Helmholtz.}
\end{figure}

\part[20]
A medi\c{c}\~{a}o do campo foi feita com recurso a uma sonda de efeito de Hall. Explique o que \'{e} o efeito de Hall num semicondutor cujos portadores de carga maiorit\'{a}rios s\~{a}o negativos. (Nota: Pode recorrer a esbo\c{c}o gr\'{a}fico para ilustrar a sua resposta.)

\part[20]
A calibra\c{c}\~{a}o da sonda foi feita recorrendo a um solenoide padr\~{a}o com $3467\pm60$ espiras por metro. As medidas obtidas est\~{a}o presentes na Tabela~\ref{tab:calibhall}.

\begin{table}[h]
\caption{\label{tab:calibhall}Tens\~{a}o ($V_H$) amplificada da sonda por efeito de Hall em fun\c{c}\~{a}o da corrente ($I_S$) que atravessa o solenoide padr\~{a}o.}
\begin{center}
\begin{tabular}{|c|c|}
\hline
$I_S~\pm~0.01$ (A) & $V_H~\pm~0.1$ (mV) \\ \hline
0.08 & 32.5 \\ 
0.18 & 48.5 \\ 
0.29 & 64.5 \\ 
0.47 & 91.9 \\ 
0.69 & 126.4 \\ 
0.82 & 145.8 \\ \hline
\end{tabular}
\end{center}
\end{table}

Com base nestes valores, obtenha a constante de calibra\c{c}\~{a}o da sonda de efeito de Hall, bem como o respetivo erro. Nota: O campo magn\'{e}tico ao longo do eixo de um solenoide infinito \'{e} dado pela Eq.~\ref{eq:solenoide}.
\begin{equation}
\label{eq:solenoide}
\left|\overrightarrow{B_S}\right|=\mu_0\frac{N}{l}I_S
\end{equation}

\part[15]
Determine a corrente aplicada numa bobina se esta tiver $200~\text{espiras}$, sabendo que o campo magn\'{e}tico produzido por uma espira ao longo do seu eixo \'{e} dado pela Eq.~\ref{eq:coilfield} e que o valor m\'{a}ximo medido da tens\~{a}o de Hall \'{e} $V_{H,\text{max}}=38.8~\text{mV}$.

\begin{equation}
\label{eq:coilfield}
B\left(x\right)=\frac{\mu_0}{2}\frac{IR^2}{\left(R^2+x^2\right)^{3/2}}
\end{equation}
\end{parts}

\end{questions}
\end{document}